Dues esferes iguals de 20~g de massa pengen cadascuna d'un fil de 50~cm de llarg, tal com mostra la figura. Totes dues esferes tenen càrregues elèctriques iguals, però de signe contrari. A causa de l'atracció elèctrica que hi ha entre les esferes, els fils formen un angle de $15^\circ$ amb la vertical. En aquesta configuració, la distància entre les esferes és de 10~cm.

\figura[0.35]{fig1.png}

\begin{apartat}{1,25}
Calculeu el mòdul de la força elèctrica entre les esferes i el valor de les seves càrregues elèctriques.
\end{apartat}

\begin{apartat}{1,25}
Si retiréssim la càrrega positiva, quin camp hauríem de crear al voltant de la càrrega negativa perquè aquesta última no canviés de posició? Indiqueu-ne el mòdul i representeu esquemàticament la direcció i el sentit que tindria. Com hauria de ser aquest camp si, en lloc de retirar la càrrega positiva, retiréssim la càrrega negativa?
\end{apartat}

\blocetiquetat{Dada:}{%
  $k = \dfrac{1}{4\pi\varepsilon_0} = 8,99\times 10^{9}\ \mathrm{N\,m^2\,C^{-2}}$.}
\nota{Suposeu que les dues càrregues són puntuals.}
